\subsubsection{Data Persistence}
\begin{enumerate}
    \item \textbf{Python Interface for Array Serialization and Deserialization}
          
          The Python interface provides bindings to the Fortran serialization and deserialization
          routines, enabling efficient handling of multi-dimensional arrays.This
          interface supports NumPy arrays with native data type compatibility and
          Fortran-style column-major memory layout.
          \begin{enumerate}
            \item \textbf{General Architecture}
                  
                  The Python module \texttt{tensoromics\_functions} provides ten main
                  functions for array handling:
                  
                  \begin{itemize}
                    \item \texttt{tox\_serialize\_int\_nd(array, filename)} - Serialize
                          N-dimensional integer arrays.
                          
                    \item \texttt{tox\_deserialize\_int\_nd(filename)} - Deserialize N-dimensional
                          integer arrays.
                          
                    \item \texttt{tox\_serialize\_real\_nd(array, filename)} - Serialize
                          N-dimensional real arrays.
                          
                    \item \texttt{tox\_deserialize\_real\_nd(filename)} - Deserialize N-dimensional
                          real arrays.
                          
                    \item \texttt{tox\_serialize\_char\_nd(array, filename)} - Serialize
                          N-dimensional character arrays.
                          
                    \item \texttt{tox\_deserialize\_char\_nd(filename)} - Deserialize N-dimensional
                          character arrays.
                          
                    \item \texttt{tox\_serialize\_logical\_nd(array, filename)} -
                          Serialize N-dimensional logical arrays.
                          
                    \item \texttt{tox\_deserialize\_logical\_nd(filename)} - Deserialize
                          N-dimensional logical arrays.
                          
                    \item \texttt{tox\_serialize\_complex\_nd(array, filename)} -
                          Serialize N-dimensional complex arrays.
                          
                    \item \texttt{tox\_deserialize\_complex\_nd(filename)} - Deserialize
                          N-dimensional complex arrays.
                          
                    \item \texttt{get\_array\_metadata(filename)} - Retrieve array
                          dimensions and metadata
                  \end{itemize}
                  
            \item \textbf{Serialization Functions:}
                  \begin{itemize}
                    \item \texttt{array} (input): NumPy array to be serialized \\ \textbf{Type:}
                          \texttt{numpy.ndarray} with dtype:
                          \begin{itemize}
                            \item \texttt{np.int32} for integer
                                  
                            \item \texttt{np.float64} for float values
                                  
                            \item \texttt{U5} for characters
                                  
                            \item \texttt{boolean} for logical
                                  
                            \item \texttt{np.complex128} for complex numbers
                          \end{itemize}
                          \textbf{Layout:} Column-major (Fortran) order (\texttt{order='F'})
                          \\ \textbf{Dimensions:} Since the array is being passed flat,
                          theoretically all dimensions are supported, but practical use cases
                          are typically 1D to 5D
                          
                    \item \texttt{filename} (input): Name of the binary file to create \\
                          \textbf{Type:} \texttt{str}
                  \end{itemize}
                  \begin{lstlisting}[language=Python, caption={Example usage of serialization function}]
tox_serialize_int_nd(int_array, filename)
tox_serialize_real_nd(real_array, filename)
                  \end{lstlisting}
                  
            \item \textbf{Deserialization Functions:}
                  \begin{itemize}
                    \item \texttt{filename} (input): Name of the binary file to read \\ \textbf{Type:}
                          \texttt{str}
                          
                    \item \textbf{Return Value:} Deserialized NumPy array \\ \textbf{Type:}
                          \texttt{numpy.ndarray} with original dtype and dimensions \\
                          \textbf{Layout:} Column-major (Fortran) order
                  \end{itemize}
                  
                  \begin{lstlisting}[language=Python, caption={Example usage of deserialization function}]
restored_int = tox_deserialize_int_nd(filename)
restored_real = tox_deserialize_real_nd(filename)
                  \end{lstlisting}
                  
            \item \textbf{Implementation Details}
                  
                  \begin{itemize}
                    \item The Python interface uses Fortran-style column-major array ordering
                          for optimal compatibility
                          
                    \item All arrays are converted to the appropriate Fortran data types
                          before serialization
                          
                    \item The interface automatically handles metadata including dimensions
                          and data types
                          
                    \item Error handling is implemented through Python exceptions
                  \end{itemize}
                  
            \item \textbf{Usage Example}
                  
                  \begin{lstlisting}[language=Python, caption={Full example of serialization and deserialization}]
# Real array example
real_array = np.array([1.5, 2.3, 3.2, 4.0, 5.0], dtype=np.float64, order='F')
tox_serialize_real_nd(real_array, "test_real_1d.bin")
restored_real = tox_deserialize_real_nd("test_real_1d.bin")

# Character array example with metadata
char_array = np.asfortranarray(["hello", "world"], dtype='U5')
tox_serialize_char_nd(char_array, "test_char_1d.bin")

restored_char = tox_deserialize_char_nd("test_char_1d.bin")

# Multi-dimensional examples
int_2d = np.array([[1, 2, 3], [4, 5, 6]], dtype=np.int32, order='F')
tox_serialize_int_nd(int_2d, "test_int_2d.bin")

restored_2d = tox_deserialize_int_nd("test_int_2d.bin")
                  \end{lstlisting}
                  
            \item \textbf{Compatibility Notes}
                  
                  \begin{itemize}
                    \item Requires NumPy installation
                          
                    \item Arrays must use Fortran-style column-major ordering
                          
                    \item Character arrays must use fixed-length string data types
                          
                    \item The interface supports N dimensions, but practical use cases
                          are typically 1D to 5D
                          
                    \item Data types are restricted to int32, float64, and U5 for characters
                  \end{itemize}
          \end{enumerate}
          
    \item \textbf{Standard Archive Operations}
          \begin{itemize}
            \item \textbf{Create Zip Archive (Low-Level)}
                  
                  The low-level function creates a zip archive from keys and filenames,
                  directly calling the Fortran backend. Use this function for custom data.
                  Note that the identifiers at position $i$ belongs to position $i$ of
                  the filenames. All files named in the filenames have to be saved to the
                  disk before the function is called. For more details on how to save
                  arrays, see section array serialization and deserialization. Call the
                  \texttt{create\_zip\_archive} function with the following arguments:
                  
                  \begin{itemize}
                    \item \texttt{zip\_filename}: Name of the zip file to create
                          
                    \item \texttt{keys}: List of keys for the manifest (identifiers for
                          each file)
                          
                    \item \texttt{filenames}: List of filenames to include in archive
                  \end{itemize}
                  
                  \begin{lstlisting}[language=Python, caption={Python Create Zip Archive (Low-Level)}]
create_zip_archive(zip_filename, keys, filenames)
                  \end{lstlisting}
                  
            \item \textbf{Extract Zip Archive}
                  
                  The function extracts a zip archive created by \texttt{create\_zip\_archive}
                  and returns a dictionary mapping data keys to extracted filenames. All
                  extracted files will be written to the disk. For more details on how
                  to continue working with the files, see section array serialization and
                  deserialization. Call the \texttt{extract\_zip\_archive} function with
                  the following arguments:
                  
                  \begin{itemize}
                    \item \texttt{zip\_filename}: Path to the zip file to extract
                  \end{itemize}
                  
                  \textbf{Returns:}
                  \begin{itemize}
                    \item Dictionary mapping data keys to extracted temporary filenames
                  \end{itemize}
                  
                  \begin{lstlisting}[language=Python, caption={Python Extract Zip Archive}]
extract_zip_archive(zip_filename)
                  \end{lstlisting}
                  
            \item \textbf{Save TOX Data (High-Level)}
                  
                  The high-level function saves TOX data to a zip archive, handling
                  serialization, and archive creation for standard-conform tox gene data.
                  This function automatically manages temporary files and cleanup. Call the
                  \texttt{save\_tox\_data} function with the following arguments:
                  
                  \begin{itemize}
                    \item \texttt{zip\_filename}: Name of the zip file to create
                          
                    \item \textbf{(Optional)} \texttt{gene\_ids}: Array of gene IDs
                          
                    \item \textbf{(Optional)} \texttt{expression\_vectors}: 2D array of
                          expression data
                          
                    \item \textbf{(Optional)} \texttt{gene\_to\_fam}: Array mapping
                          genes to families
                          
                    \item \textbf{(Optional)} \texttt{family\_ids}: Array of family IDs
                          
                    \item \textbf{(Optional)} \texttt{family\_centroids}: 2D array of
                          family centroids
                          
                    \item \textbf{(Optional)} \texttt{shift\_vectors}: 2D array of shift
                          vectors
                          
                    \item Corresponding \texttt{*\_name} parameters for temporary
                          filenames (e.g., \texttt{gene\_ids\_name}, \texttt{expression\_vectors\_name},
                          etc.)
                  \end{itemize}
                  
                  \begin{lstlisting}[language=Python, caption={Python Save TOX Data}]
save_tox_data(zip_filename, gene_ids=None, expression_vectors=None, 
        gene_to_fam=None, family_ids=None, family_centroids=None, 
        shift_vectors=None, gene_ids_name=None, 
        expression_vectors_name=None, gene_to_fam_name=None, 
        family_ids_name=None, family_centroids_name=None, 
        shift_vectors_name=None)
                  \end{lstlisting}
                  
            \item \textbf{Read TOX Data}
                  
                  The function reads data from a zip archive created by \texttt{save\_tox\_data},
                  extracting and deserializing the requested data components for standard-conform
                  tox gene data. This function automatically handles file extraction, cleanup
                  and array allocation. Call the \texttt{read\_tox\_data} function with the
                  following arguments:
                  
                  \begin{itemize}
                    \item \texttt{zip\_filename}: Name of the zip file to read
                          
                    \item \textbf{(Optional)} \texttt{load\_gene\_ids}: Whether to load
                          gene IDs (default: False)
                          
                    \item \textbf{(Optional)} \texttt{load\_expression\_vectors}:
                          Whether to load expression vectors (default: False)
                          
                    \item \textbf{(Optional)} \texttt{load\_gene\_to\_fam}: Whether to
                          load gene-to-family mapping (default: False)
                          
                    \item \textbf{(Optional)} \texttt{load\_family\_ids}: Whether to
                          load family IDs (default: False)
                          
                    \item \textbf{(Optional)} \texttt{load\_family\_centroids}: Whether
                          to load family centroids (default: False)
                          
                    \item \textbf{(Optional)} \texttt{load\_shift\_vectors}: Whether to
                          load shift vectors (default: False)
                  \end{itemize}
                  
                  \textbf{Return dictionary:}
                  \begin{itemize}
                    \item \textbf{gene\_ids}: Loaded gene IDs array (if requested)
                          
                    \item \textbf{expression\_vectors}: Loaded expression vectors array (if
                          requested)
                          
                    \item \textbf{gene\_to\_fam}: Loaded gene-to-family mapping array (if
                          requested)
                          
                    \item \textbf{family\_ids}: Loaded family IDs array (if requested)
                          
                    \item \textbf{family\_centroids}: Loaded family centroids array (if
                          requested)
                          
                    \item \textbf{shift\_vectors}: Loaded shift vectors array (if
                          requested)
                  \end{itemize}
                  
                  \begin{lstlisting}[language=Python, caption={Python Read TOX Data}]
read_tox_data(zip_filename, load_gene_ids=False, load_expression_vectors=False, 
            load_gene_to_fam=False, load_family_ids=False, 
            load_family_centroids=False, load_shift_vectors=False)
                  \end{lstlisting}
                  
                  \textbf{Data Archiving Workflow Examples:}
                  \begin{lstlisting}[language=Python, caption={Python Standard TOX Data Archiving}]
# Save complete TOX dataset
save_tox_data(
    zip_filename="complete_tox_dataset.zip",
    gene_ids=filtered_gene_ids,
    gene_ids_name="gene_ids.bin",
    expression_vectors=filtered_expression,
    expression_vectors_name="expression.bin",
    gene_to_fam=filtered_gene_to_fam,
    gene_to_fam_name="gene_to_fam.bin",
    family_ids=family_ids,
    family_ids_name="family_ids.bin",
    family_centroids=family_centroids,
    family_centroids_name="centroids.bin",
    shift_vectors=shift_vectors,
    shift_vectors_name="shift_vectors.bin"
)

# Save partial dataset
save_tox_data(
    zip_filename="basic_data.zip",
    gene_ids=filtered_gene_ids,
    gene_ids_name="gene_ids.bin",
    expression_vectors=filtered_expression,
    expression_vectors_name="expression.bin"
)

# Read data back with selective loading
loaded_data = read_tox_data(
    zip_filename="complete_tox_dataset.zip",
    load_gene_ids=True,
    load_expression_vectors=True,
    load_gene_to_fam=True
)
                  \end{lstlisting}
                  
                  \begin{lstlisting}[language=Python, caption={Python Custom Archive Creation and Extraction}]
# Create custom archive with non-standard data
create_zip_archive(
    zip_filename="custom_data.zip",
    keys=["3d_tensor", "results"],
    filenames=["tensor_data.bin", "analysis_results.bin"]
)

# Extract and access custom archive
file_mapping = extract_zip_archive("custom_data.zip")
print("Files in archive:", file_mapping)
# Output: {'3d_tensor': 'tensor_data.bin', 'results': 'analysis_results.bin'}
                  \end{lstlisting}
                  
                  \begin{lstlisting}[language=Python, caption={Python Mixed Standard and Custom Data Archiving}]
# Serialize custom arrays
tox_serialize_real_nd(custom_array, "custom_data.bin")
tox_serialize_int_nd(another_array, "another_data.bin")

# Create mixed archive with both standard and custom data
create_zip_archive(
    zip_filename="mixed_archive.zip",
    keys=["standard_expression", "custom_feature", 
            "analysis_results"],
    filenames=["expression.bin", "custom_data.bin", 
            "another_data.bin"]
)

# Later, extract and process
mapping = extract_zip_archive("mixed_archive.zip")
expression_data = tox_deserialize_real_nd(
    mapping["standard_expression"])

custom_features = tox_deserialize_real_nd(
    mapping["custom_feature"])

results = tox_deserialize_int_nd(mapping["analysis_results"])
                  \end{lstlisting}
          \end{itemize}
\end{enumerate}
